% Record of research-paper questions and answers from the conversation on
% 2026-11-16. Answers are traced to specific files/functions in the NATM
% codebase as of that date; re-verify against the code if it changes later.
% No LaTeX compiler was available in the environment that produced this
% file, so it has NOT been compiled/verified -- please compile it yourself
% and report back if anything breaks.

\documentclass{article}
\usepackage[margin=1in]{geometry}
\usepackage{enumitem}
\usepackage{array}
\usepackage{longtable}
\usepackage{hyperref}

\title{NATM Framework --- Research Questions and Answers}
\date{2026-11-16}

\begin{document}
\maketitle

\section*{Questions}

\begin{enumerate}[leftmargin=*]
  \item When can an aircraft or vessel be replaced? Is replacement triggered
    by age or technical lifetime? Can operators replace or retrofit assets
    early? If transport activity grows, how does the model decide whether to
    use existing assets more intensively or add new ones?

  \item Which technologies and fuels were actually available? Please give
    the options enabled for passenger aviation, air cargo and maritime
    transport. What makes an option feasible for a specific asset or route,
    such as range, vessel class, commercial availability or access to fuel?
    Were fuel mandates applied as constraints or reflected in costs?

  \item What enters the NPV, and what varies across the six Ariadne
    pathways? Which costs and revenues are included for aircraft and
    vessels? In particular, which of the electricity, fuel and carbon
    prices, activity projections and policy assumptions vary by pathway, and
    which are shared?

  \item How does the 2025 activity baseline grow into future route
    activity? The baseline comes from observed flights and sailings. Are
    Ariadne growth rates applied uniformly to those routes, or do they
    differ by segment, route or vessel class? Can fleet renewal change which
    routes an asset serves?

  \item How are fleet choices converted into the demand profile passed to
    the market assessment? What determines energy use per flight or voyage,
    how is it attributed to airports or ports, and which output becomes the
    single profile used across all six market pathways? If the fleet model
    produces pathway-specific operating demand, which trajectory is passed
    onward?
\end{enumerate}

\section*{Answers}

Every claim below is traced to a specific file or function in the
codebase, not reconstructed from memory. Where something is a documented
limitation rather than a finding, that is stated explicitly.

\subsection*{1. Replacement: age-triggered, with a narrow early-replacement
window; growth favors existing assets first}

\textbf{The trigger is technical lifetime, not calendar age.} On creation,
each asset is given \verb|replacement_year = start_year + (lifetime_years|
\verb|- current_age)| (\texttt{fleet.py::\_initial\_replacement\_year}); after
any technology swap, \verb|replacement_year = year + lifetime_years| of the
newly chosen technology (\texttt{fleet.py::apply\_technology}, line 475).
Replacement becomes ``due'' purely when
\verb|replacement_year <= current_year|
(\texttt{fleet.py::due\_replacement\_indices}, line 442).

\textbf{Early replacement is possible, but narrowly.} It is available only
for assets still on a \emph{conventional} technology, only within an
\verb|acceleration_window| computed from policy pressure ---
\verb|acceleration_window = int((adoption_mandate + clean_fuel_subsidy) * 5)|
(\texttt{scorer.py::replacement\_row\_indices}) --- and only once a
\verb|minimum_holding_period| has elapsed since the asset's last investment.
Retrofitting is therefore driven by policy/subsidy strength, not operator
discretion, and is bounded by a minimum ownership period.

\textbf{Growth favors intensifying existing assets before adding new
ones, and this is capped, not elastic.} Each year,
\verb|segment_utilization_factor = min(1.0, allocated_demand / full_capacity)|
scales how much of the \emph{existing} fleet's real baseline distance is
actually flown (\texttt{aviation\_passenger\_airline.py}, line 269, mirrored
in the cargo and maritime agents) --- demand below full fleet capacity is
absorbed by flying the existing fleet \emph{less}, never more than its real
base-year rate. Only when demand exceeds full existing capacity does
\verb|_add_growth_assets| (\texttt{scorer.py}, line 180) add new assets,
sized to close exactly that gap.

This utilization-factor mechanism has a documented history worth citing
directly: the code comment records that before it existed, ``every aircraft
always flew its fixed, real \verb|annual_distance_km_base| every year
regardless of\ldots scenario,'' which the comment states was
\emph{confirmed} to make total passenger energy demand byte-identical
across all 8 scenarios despite demand differing by up to $\sim$19\%. That is
a concrete, verifiable before/after in the model's own development history,
worth citing if the paper discusses model validation.

\subsection*{2. Technology and fuel options, feasibility, and mandates}

Full enabled catalog, current state (service-entry year in parentheses):

\begin{longtable}{p{2.6cm}p{3.4cm}p{3.4cm}p{5.6cm}}
\textbf{Sector} & \textbf{Conventional} & \textbf{Bio/drop-in} & \textbf{Synthetic/alternative} \\
\hline
\endhead
Passenger aviation &
  kerosene short/medium/long (1995--2000); efficient turbofan short/medium (2035) &
  drop-in SAF (HEFA) short/medium/long (2025) &
  e-kerosene (PtL) short/medium/long (2028); hydrogen short (2035); hybrid-electric short (2031) \\
Cargo aviation &
  kerosene freight medium/long (2005--2008) &
  drop-in SAF (HEFA) freight medium/long (2025) &
  e-kerosene freight medium/long (2028); hydrogen freight medium (2030) \\
Maritime cargo &
  marine diesel regional/deepsea (2010--2011) &
  biofuel (FAME) regional/deepsea (2025) &
  LNG (2015); methanol (2023); ammonia (2030--2031); hydrogen regional (2032); battery regional (2022); hybrid regional/deepsea (2018); wind-assist (2015); plus ``ready'' (retrofit-prepared, still-diesel) variants of methanol/ammonia \\
\end{longtable}

\textbf{Feasibility is three independent gates}, applied in this order
(\texttt{technology\_catalog.py::candidates\_for\_operation}, then
\texttt{legacy\_weighted\_utility.py::is\_candidate\_available}):
\begin{enumerate}
  \item \textbf{Segment and range}: a candidate's \verb|trip_length_km| must
    meet the asset's required stage length for its segment; if \emph{no}
    candidate in the catalog clears that bar, the filter is skipped entirely
    rather than leaving zero candidates.
  \item \textbf{Service-entry year}: a hard cutoff --- unavailable before its
    \verb|service_entry_year|, no exception.
  \item \textbf{Scenario-driven availability flags}: \verb|technology_availability|,
    \verb|infrastructure_availability|, and (only for blended/drop-in fuels)
    \verb|saf_availability| --- each defaults to available (1.0) when the
    scenario file carries no explicit row for that technology/segment/year.
\end{enumerate}

\textbf{One point worth flagging precisely}: the catalog carries
\verb|minimum_airport_class|/\verb|minimum_port_class| columns, but these are
confirmed to be \textbf{never read anywhere in the decision logic} ---
descriptive metadata only, not an active constraint. If the paper references
vessel-class or airport-class feasibility, that gate does not currently
exist in the code, only in the data schema.

\textbf{Mandates are a hard physical constraint, not a cost nudge.}
\verb|effective_secondary_share| sets the actual blend fraction burned to
\verb|min(technology_max_share, mandate_share)| once a mandate is active
(\texttt{legacy\_weighted\_utility.py::effective\_secondary\_share}, line
120) --- this directly sets \verb|primary_energy_quantity| and
\verb|secondary_energy_quantity| in \verb|annual_operation_metrics|, which is
what cost and emissions are computed from. The mandate does not compete for
selection via utility scoring; it mechanically forces the fuel split, and
cost follows from that. (A separate, smaller \verb|policy_bonus| term exists
in the \emph{legacy} weighted-utility scorer, but it is not part of the NPV
figure the ambiguity-aware selection --- which is what every one of the 24
simulated cases actually uses --- optimizes on; see Question 3.)

\subsection*{3. What enters NPV, and the pathway-variation table}

NPV (\texttt{scorer.py::calc\_payback\_year}) equals $-$CAPEX at decision
time, plus the discounted sum over the technology's lifetime of (revenue
$-$ operating cost $-$ additional costs $-$ depreciation), plus a discounted
residual value at the end. Concretely:
\begin{itemize}
  \item \textbf{Revenue}: payload/seat capacity $\times$ load factor
    $\times$ distance flown $\times$ fare/freight rate (sector-specific
    \verb|annual_revenue|).
  \item \textbf{Operating cost}: \verb|energy_cost| (primary $+$ secondary
    fuel quantity $\times$ scenario price) $+$ \verb|emission_cost|
    (chargeable emissions $\times$ carbon price, net of free ETS
    allocation).
  \item \textbf{Additional costs}: maintenance
    (\verb|maintenance_cost_share| $\times$ revenue) plus fixed operational
    overheads as a percentage of revenue --- crew/port-fees/cargo-handling
    for maritime, wages/landing-fees for aviation.
  \item \textbf{Depreciation}: \verb|depreciation_cost_share| $\times$
    CAPEX, charged every year of the holding period.
\end{itemize}

Economic/environmental ``utility'' and the policy bonus are computed and
recorded for diagnostics, but are \textbf{not part of the NPV figure
itself} --- confirmed directly in
\texttt{ambiguity\_aware\_utility.py::\_decision\_score}, whose comment
states the ambiguity-aware selection criterion is net present value, ``not
the normalized payback-year heuristic.''

Pathway variation, verified directly across all three sectors' scenario
files --- an identical pattern in every one:

\begin{longtable}{p{7.5cm}p{7.5cm}}
\textbf{Varies by pathway (6 distinct values)} & \textbf{Shared across all 6 pathways} \\
\hline
\endhead
Primary energy price, secondary energy price (maritime also: tertiary energy price) &
  Carbon price/tax, clean fuel subsidy, ETS allocation factor, freight rate, load factor, technology/infrastructure availability, operator market share, planned delivery count \\
Demand (passenger-km / freight-tonne-km) & --- \\
SAF/biofuel mandate (aviation only --- maritime's \texttt{biofuel\_mandate} is shared, not pathway-varying) & --- \\
\end{longtable}

\textbf{One caveat to state precisely in the paper}: cargo aviation's
\verb|freight_tonne_km_demand| technically has six separate rows (one per
pathway) but --- confirmed directly, and confirmed as an intentional design
choice rather than a data bug --- the six rows carry \textbf{identical
values}. For cargo aviation specifically, demand does not actually
differentiate the six pathways even though the data schema allows it to;
passenger aviation's and maritime's demand do genuinely differ by pathway.

\subsection*{4. Baseline growth: segment-level, not route-level; fleet
renewal never reassigns routes}

\textbf{Growth rates differ by segment, not by individual route or vessel
class.} Demand is scenario-scoped at country $\times$ segment only
(\texttt{short}/\texttt{medium}/\texttt{long} for aviation,
\texttt{regional}/\texttt{deepsea} for maritime) --- confirmed directly
against the scenario files; there is no per-route or per-operator demand
row. Within a segment, growth is uniform across every route/asset in that
segment; across segments, trajectories differ (e.g.\ passenger
short/medium/long diverge by several percentage points by 2050 under the
same pathway).

\textbf{Fleet renewal never changes which route or hub an asset serves.}
\verb|apply_technology| (\texttt{fleet.py}, line 457) updates technology,
carrier, replacement year, age, and financial/emissions fields --- it does
\textbf{not} touch \verb|main_hub_base|, \verb|segment|, or any
route-identifying field. A technology swap changes \emph{what} an asset
burns, never \emph{where} it flies. Growth additions are seeded from the
same segment's template, so even new aircraft inherit that segment's
characteristic route profile rather than being freely routable.

\subsection*{5. Fleet $\to$ demand profile $\to$ market assessment: one
important limitation to flag}

Energy use per flight/voyage is \verb|annual_operation_metrics|'s
\[
\text{total\_energy} = \frac{\text{real baseline distance} \times \text{utilization\_factor}}{\text{technology efficiency}},
\]
split into primary/secondary by the (mandate-determined) blend share. This
is attributed to a single real hub per asset (\verb|main_hub_base|) --- not
a multi-airport route network; that is the basis for every per-airport
total in the dashboard and the 15-minute export. Energy per route (not just
per hub) \emph{is} separately available via the OpenAP-derived route
summary, but that is an external post-processing join, not something the
simulation itself attributes.

\textbf{On ``which single profile gets passed to the market
assessment''}: precision over speculation here. A search of this
repository's documentation found no ``market assessment'' component, no
Ariadne-coupling description, and no aggregation step that merges
pathway-specific outputs into one. What exists is the
\verb|realized_scenario_id| mechanism: each \emph{case} (e.g.\
\texttt{fokusH2-risk-neutral}) is configured with one
\verb|realized_scenario_id|, and that case's \texttt{aircraft.csv} /
\texttt{*\_energy\_emissions.csv} contains the single simulated trajectory
under that one realized pathway --- not six parallel trajectories merged
into one. There is no NATM-internal step that picks ``the'' profile out of
six; \textbf{there are six (or eight, counting the two risk-averse cases)
separate single-pathway output sets}, one per case, and whichever one feeds
``the market assessment'' is a choice made outside this codebase --- most
likely by running the case whose \verb|realized_scenario_id| matches the
pathway the market assessment is evaluating. Stating precisely how that
handoff happens would describe an external process not visible from this
repository alone.

\end{document}
